\chapter{Introduction}
\label{ch:intro}

\section{Context and objective}

Traction inverters are the point where the electrical and the mechanical
domains of an electric vehicle meet, and they are the component in which the
move from silicon to wide-bandgap devices has had the clearest effect. Raising
the DC-link from the \SI{400}{\volt} class to the \SI{800}{\volt} class halves
the current for a given power, which reduces conduction losses and cable
cross-section, but it also puts the switching devices in a voltage class where
silicon IGBTs become slow and lossy. Silicon carbide MOSFETs are the natural
answer, and they bring their own set of problems: fast edges, small gate-drive
margins, and a body diode whose reverse recovery is not negligible.

The object of this work is the power stage of such a converter: the bridge,
the DC-link and the cooling, with the schematic and the board layout of that
module. \Cref{fig:system} shows what the module contains and where it stops:
the battery, the gate drive and the controller are outside it, and what
crosses its boundary is written down as an assumption (\cref{tab:sysassume}).
The report follows the four things the brief asks of it: the design
equations (\cref{ch:analytical,ch:modulation}), the component selection
(\cref{ch:spec}), the design choices, and the discussion of the numerical
simulations (\cref{ch:model,ch:losses,ch:thermal}). Every component value is
either taken from a datasheet, computed from a formula shown in the text, or
obtained from a simulation whose model is documented and validated.

\begin{figure}[htbp]
\centering
% Il power stage e i suoi confini: cosa sta dentro il modulo e cosa gli arriva da fuori.
\begin{tikzpicture}[
  font=\footnotesize, >=Stealth,
  blk/.style={draw, rounded corners=2pt, align=center, minimum height=10mm, inner sep=3pt, fill=white},
  hv/.style={blk, draw=farancio, fill=farancio!7},
  ext/.style={blk, dashed, draw=fgrigio},
  pw/.style={->, very thick, farancio}]
  % catena di potenza
  \node[ext, text width=17mm] (bat) at (0,0) {HV battery\\\SI{800}{\volt}};
  \node[hv, text width=26mm] (dcl) at (4.0,0) {DC-link\\film and ceramics};
  \node[hv, text width=38mm] (brg) at (8.6,0) {three-phase bridge\\$6\times$ SiC MOSFET, shunts};
  \node[draw, circle, minimum size=11mm] (mot) at (13.0,0) {M};
  \draw[pw] (bat) -- node[above, font=\scriptsize]{DC} (dcl);
  \draw[pw] (dcl) -- (brg);
  \draw[pw] (brg) -- node[above, font=\scriptsize]{3 phases} (mot);
  % raffreddamento: dentro il modulo, con il suo liquido fuori
  \node[hv, text width=38mm] (rail) at (8.6,1.9) {cooling rail, pad\\{\scriptsize liquid at \SI{65}{\celsius}}};
  \draw[->, thick, fgrigio] (brg) -- node[right, font=\scriptsize]{heat} (rail);
  % cio' che arriva da fuori
  \node[ext, text width=46mm] (ctl) at (8.6,-2.1) {gate drive, protection, controller\\{\scriptsize outside the module}};
  \draw[->, thick] (ctl.north) -- node[right, font=\scriptsize, pos=0.28]{gate and Kelvin-source signals} (brg.south);
  % confine del modulo
  \begin{scope}[on background layer]
    \node[draw=black, dash dot, rounded corners=4pt, inner sep=3mm,
          fit=(dcl)(brg)(rail)] (mod) {};
  \end{scope}
  \node[anchor=north west, font=\scriptsize\bfseries] at (mod.north west) {power stage module};
\end{tikzpicture}

\caption{The power stage and its boundaries. Orange: inside the module; dashed
grey: outside it, with the interface stated as assumptions
(\cref{tab:sysassume}).}
\label{fig:system}
\end{figure}

\section{Method and structure}

The design proceeds in four stages, each validating the previous one
(\cref{fig:flow}). Analytical sizing fixes the operating point, the DC-link
capacitance, the device rating and the modulation from closed-form relations
(\cref{ch:analytical,ch:modulation}). Circuit simulation builds a PLECS model
of the converter and uses it to obtain the losses; its outputs are reconciled
with the analytical results, the vendor loss tables it relies on are verified
against the datasheet, and a SPICE study of a single switching cell fixes the
gate resistances on the overvoltage at the die (\cref{ch:model,ch:losses}).
Thermal analysis closes the loop between losses and temperature, which this
device needs because its on-resistance rises by \SI{80}{\percent} between
\SI{25}{\celsius} and \SI{175}{\celsius}, and verifies the stalled-motor case
that selected the device (\cref{ch:thermal}). Schematic and layout turn the
design into a KiCad module (\cref{ch:schematic,ch:layout}).
\Cref{ch:conclusions} draws the conclusions and \cref{app:assumptions}
collects the assumptions that still need confirmation.

\begin{figure}[htbp]
\centering
% Il metodo: le fasi in avanti e i ritorni che hanno cambiato il progetto.
\begin{tikzpicture}[
  font=\footnotesize, >=Stealth,
  fase/.style={draw=fblu, fill=fblu!7, rounded corners=2pt, align=center, text width=18mm,
               minimum height=14mm, inner sep=2pt},
  ritorno/.style={->, farancio, thick, rounded corners=3pt},
  lab/.style={font=\scriptsize, farancio, fill=white, inner sep=1.5pt}]
  \node[fase] (s) at (0,0)     {specification\\and criteria\\[1pt]\scriptsize Ch.~\ref{ch:spec}};
  \node[fase] (a) at (2.2,0)   {analytical\\sizing\\[1pt]\scriptsize Ch.~\ref{ch:analytical}--\ref{ch:modulation}};
  \node[fase] (m) at (4.4,0)   {circuit models\\{\scriptsize PLECS, SPICE}\\[1pt]\scriptsize Ch.~\ref{ch:model}--\ref{ch:losses}};
  \node[fase] (t) at (6.6,0)   {thermal\\and cooling\\[1pt]\scriptsize Ch.~\ref{ch:thermal}};
  \node[fase] (x) at (8.8,0)   {auxiliary\\circuits\\[1pt]\scriptsize Ch.~\ref{ch:aux}};
  \node[fase] (b) at (11.0,0)  {schematic\\and board\\[1pt]\scriptsize Ch.~\ref{ch:schematic}--\ref{ch:layout}};
  \node[fase, draw=fgrigio, fill=fgrigio!10] (v) at (13.2,0) {bench plan\\[1pt]\scriptsize App.~\ref{app:assumptions}};
  \foreach \i/\j in {s/a, a/m, m/t, t/x, x/b, b/v} \draw[->, thick, fblu] (\i) -- (\j);
  % ritorni: ciascuno e' una scoperta che ha cambiato una decisione a monte
  \draw[ritorno] (6.6,-0.7) -- (6.6,-1.15) -- node[lab]{locked rotor changes the device} (0,-1.15) -- (0,-0.7);
  \draw[ritorno] (11.0,-0.7) -- (11.0,-1.75) -- node[lab]{drawn rail replaces the ideal plate} (6.9,-1.75) -- (6.9,-0.7);
  \draw[ritorno] (4.1,0.7) -- (4.1,1.15) -- node[lab]{C2 sets $R_g$ and dead time} (0,1.15) -- (0,0.7);
  \draw[ritorno] (11.0,0.7) -- (11.0,1.75) -- node[lab]{extracted loop re-simulated against C2} (4.7,1.75) -- (4.7,0.7);
\end{tikzpicture}

\caption{Design flow. Blue: the stages in their chapters; orange: the four returns
that changed an earlier decision: the locked rotor that changed the
device, the overvoltage criterion that set the gate resistances and the
dead time, the extracted commutation loop simulated again, and the drawn
cooling rail that replaced the ideal cold plate.}
\label{fig:flow}
\end{figure}

Four times a later stage contradicted an assumption made earlier, and each
return is reported where it happened, in a shaded box. The ones that changed
the design are the choice of device (\cref{sec:deviceselection}), the
double-pulse bench that measured a charging bus (\cref{sec:kelvin}), the
interpretation of the vendor loss library (\cref{sec:losstables}), the
dead-time loss calculation (\cref{sec:deadtime-loss}), the identification of
the dominant thermal unknown (\cref{sec:rthcs}) and the cooling rail, whose
real resistance an ideal cold plate had hidden (\cref{sec:cooling}).
